\section{Finite Selberg estimates with explicit constants}
\label{sb:section}

We now prove all the number-theoretic estimates needed for the contraction
argument.  The constants are deliberately generous.  Their purpose is to
make every error bound visible and to avoid any appeal to an unspecified
asymptotic estimate.

For a positive integer $n$, define
\[
 \Lambda(n)=
 \begin{cases}
  \log p,&n=p^j\text{ for a prime }p\text{ and }j\ge1,\\
  0,&\text{otherwise},
 \end{cases}
 \qquad
 \psi(x)=\sum_{n\le x}\Lambda(n),
 \qquad R(x)=\psi(x)-x.
\]
Thus $\Lambda(1)=0$, and $\psi(x)=0$ for $0\le x<2$.  A summation
with a real upper endpoint means a finite sum over the positive integers
below that endpoint.  All integrals below have bounded endpoints.  They
can be evaluated or approximated by the finite arithmetic procedures
specified in the arithmetization part of the paper.

\subsection{Harmonic estimates without an infinite constant}
\label{sb:harmonic}

Write
\[
 H_m=\sum_{j=1}^m\frac1j,\qquad
 h_m=H_m-\log m,\qquad
 P_i(y)=\sum_{n\le y}(\log n)^i\quad(i=1,2).
\]
We use the elementary inequalities
\begin{equation}
 \frac{z-1}{z}\le\log z\le z-1\qquad(z\ge1).
 \label{sb:log-basic}
\end{equation}
They follow by comparing $1/t$ with its endpoint values on $[1,z]$.
In particular $1/2\le\log2\le1$.

\begin{lemma}
\label{sb:harmonic-finite}
For integers $1\le m\le M$,
\[
 0\le h_m-h_M\le\frac1m-\frac1M,
 \qquad 0<h_M\le1.
\]
For every $y\ge1$,
\begin{align}
 |P_1(y)-(y\log y-y)|&\le\log y+1\le y,
 \label{sb:log-sum}\\
 |P_2(y)-y((\log y)^2-2\log y+2)|
   &\le(\log y)^2+2.
 \label{sb:log-square-sum}
\end{align}
\end{lemma}

\begin{proof}
The difference of successive $h$'s is
\[
 h_n-h_{n+1}=\log(1+1/n)-\frac1{n+1}.
\]
By \eqref{sb:log-basic}, this lies between $0$ and
$1/n-1/(n+1)$.  A finite telescoping sum proves the first assertion.
The rectangle bounds
\[
 \log(M+1)\le H_M\le1+\log M
\]
give the second assertion.

If $f$ is either $\log t$ or $(\log t)^2$, and $N=\lfloor y\rfloor$,
monotonicity on $[1,y]$ gives
\[
 0\le\sum_{n=1}^N f(n)-\int_1^N f(t)\,dt\le f(N).
\]
The extra integral from $N$ to $y$ lies between $0$ and $f(y)$.
Consequently the distance between the sum and $\int_1^y f(t)\,dt$
is at most $f(y)$.  Now use
\[
 \int_1^y\log t\,dt=y\log y-y+1,
\]
\[
 \int_1^y(\log t)^2\,dt
    =y((\log y)^2-2\log y+2)-2.
\]
Finally, $\log y+1\le y$ is another instance of
\eqref{sb:log-basic}.
\end{proof}

\subsection{A Chebyshev bound and the signed logarithmic mean}
\label{sb:chebyshev}

\begin{lemma}
\label{sb:basic-psi}
For $x\ge1$,
\begin{equation}
 0\le\psi(x)\le4x,\qquad |R(x)|\le4x,
 \label{sb:chebyshev-bound}
\end{equation}
and
\begin{equation}
 \left|\sum_{n\le x}\frac{\Lambda(n)}n-\log x\right|\le6.
 \label{sb:mertens-bound}
\end{equation}
For $1\le a\le b$,
\begin{equation}
 \left|\int_a^b\frac{R(u)}{u^2}\,du\right|\le20.
 \label{sb:signed-integral}
\end{equation}
\end{lemma}

\begin{proof}
Counting the multiples of each prime power in a factorial gives the exact
identity
\[
 \log\binom{2n}{n}
  =\sum_{d\le2n}\Lambda(d)
       \left(\left\lfloor\frac{2n}{d}\right\rfloor
              -2\left\lfloor\frac nd\right\rfloor\right).
\]
Every parenthesis is either $0$ or $1$, and it is $1$ when $n<d\le2n$.
The binomial theorem gives $\binom{2n}{n}\le2^{2n}$.  Therefore
\[
 \psi(2n)-\psi(n)\le2n\log2.
\]
Apply this for $n=1,2,\ldots,2^{k-1}$ along the dyadic sequence and
telescope:
\[
 \psi(2^k)\le(2^{k+1}-2)\log2<2^{k+1}.
\]
For $x>1$, the least dyadic integer $2^k\ge x$ is less than $2x$.
Monotonicity gives $\psi(x)\le4x$.  The case $x=1$ is immediate.
Since $R(x)=\psi(x)-x$, this also gives the stated bound on $R$.
For later use, $|R(x)|\le4x$ holds for $0\le x<1$ as well, because
there $R(x)=-x$.

The prime factorization of $n$ gives
\[
 \sum_{d\mid n}\Lambda(d)=\log n.
\]
Summing over $n\le x$ and reversing two finite sums yields
\[
 P_1(x)=\sum_{d\le x}\Lambda(d)
                      \left\lfloor\frac xd\right\rfloor.
\]
Put $A_1(x)=\sum_{d\le x}\Lambda(d)/d$.  The floor errors satisfy
\[
 |xA_1(x)-P_1(x)|\le\psi(x)\le4x.
\]
Together with \eqref{sb:log-sum}, this gives
\[
 |A_1(x)-\log x|
 \le4+\left|-1+\frac{P_1(x)-(x\log x-x)}x\right|
 \le6.
\]

The remaining identity can be checked by integrating each of the finitely
many terms $\Lambda(n)\mathbf1_{[n,x]}(t)$:
\[
 \int_1^x\frac{\psi(t)}{t^2}\,dt
     =\sum_{n\le x}\Lambda(n)\left(\frac1n-\frac1x\right)
     =A_1(x)-\frac{\psi(x)}x.
\]
Subtracting $\int_1^x dt/t=\log x$ gives
\[
 \int_1^x\frac{R(t)}{t^2}\,dt
     =A_1(x)-\frac{\psi(x)}x-\log x.
\]
Its absolute value is at most $10$.  Subtract the expressions at $b$
and $a$ to obtain \eqref{sb:signed-integral}.
\end{proof}

\subsection{A finite hyperbola calculation}
\label{sb:hyperbola}

Let $\tau(n)$ be the number of positive divisors of $n$, and put
$D(y)=\sum_{n\le y}\tau(n)$.  Fix an outer endpoint $x\ge1$ and set
\begin{equation}
 M=\lceil x\rceil,\qquad c=H_M-\log M.
 \label{sb:finite-c}
\end{equation}
The same parameter $c$ will be used for every $1\le y\le x$.
It is a finite expression; no existence or convergence assertion about
Euler's constant is being used.

\begin{lemma}
\label{sb:divisor-approx}
For $1\le y\le x$,
\begin{equation}
 |D(y)-y\log y-(2c-1)y|\le12\sqrt y.
 \label{sb:divisor-bound}
\end{equation}
Furthermore, the arithmetic function
\begin{equation}
 f(n)=\tau(n)\log n-2c\tau(n)+4c^2
 \label{sb:comparison-function}
\end{equation}
satisfies
\begin{equation}
 \left|\sum_{n\le y}\bigl(f(n)-(\log n)^2\bigr)\right|
       \le80\sqrt y(\log y+1).
 \label{sb:comparison-error}
\end{equation}
\end{lemma}

\begin{proof}
Let $m=\lfloor\sqrt y\rfloor$.  Counting pairs of positive integers
whose product is at most $y$, according to which coordinate is at most
$m$, gives
\[
 D(y)=2\sum_{n\le m}\left\lfloor\frac yn\right\rfloor-m^2
     =2yH_m-m^2+e,\qquad |e|\le2m.
\]
By Lemma~\ref{sb:harmonic-finite},
\[
 H_m=\log m+c+\eta_m,\qquad 0\le\eta_m\le1/m.
\]
Thus
\[
 D(y)-y\log y-(2c-1)y
   =-y\log(y/m^2)+(y-m^2)+2y\eta_m+e.
\]
Because $m\le\sqrt y<m+1\le2m$, we have
\[
 y\log(y/m^2)\le\frac{2y}{m}\le4\sqrt y,
 \qquad 0\le y-m^2\le2\sqrt y,
\]
\[
 2y\eta_m\le4\sqrt y,\qquad |e|\le2\sqrt y.
\]
Adding these bounds proves \eqref{sb:divisor-bound}.

For $T(y)=\sum_{n\le y}\tau(n)\log n$, finite partial summation gives
\[
 T(y)=D(y)\log y-\int_1^y\frac{D(t)}t\,dt.
\]
Write $D(t)=t\log t+(2c-1)t+E(t)$, where
$|E(t)|\le12\sqrt t$.  Substitution and elementary integration yield
\[
 T(y)=y\bigl((\log y)^2+(2c-2)\log y+2-2c\bigr)+E_T(y),
\]
where
\[
 |E_T(y)|\le2+12\sqrt y\log y+24(\sqrt y-1)
                 \le30\sqrt y(\log y+1).
\]
Here the constant term before the error is $2c-2$, of absolute value
at most $2$, since $0<c\le1$.

Now sum \eqref{sb:comparison-function}.  The coefficient of
$y\log y$ is $(2c-2)-2c=-2$, and its constant coefficient is
\[
 (2-2c)-2c(2c-1)+4c^2=2.
\]
The error in replacing $4c^2\lfloor y\rfloor$ by $4c^2y$ is at most
$4$.  It follows that
\begin{align*}
 \left|\sum_{n\le y}f(n)
           -y\bigl((\log y)^2-2\log y+2\bigr)\right|
 \le30\sqrt y(\log y+1)+24\sqrt y+4.
\end{align*}
To compare this with $P_2(y)$, use \eqref{sb:log-square-sum} and
\[
 \log y\le\int_1^y t^{-3/4}\,dt
            =4(y^{1/4}-1)\le4y^{1/4}.
\]
In particular $(\log y)^2\le16\sqrt y$.  The total error is at most
\[
 30\sqrt y(\log y+1)+40\sqrt y+6
       \le80\sqrt y(\log y+1),
\]
as required.
\end{proof}

\subsection{Selberg's symmetry formula}
\label{sb:symmetry-subsection}

For arithmetic functions on positive integers write
\[
 (a*b)(n)=\sum_{d\mid n}a(d)b(n/d).
\]
Let $\mathbf1(n)=1$ and let $\varepsilon(1)=1$,
$\varepsilon(n)=0$ for $n>1$.  The M\"obius function $\mu$ is $0$
when a square of a prime divides $n$, and otherwise is $(-1)^r$,
where $r$ is the number of distinct prime factors of $n$.
Summing over subsets of these prime factors proves
$\mu*\mathbf1=\varepsilon$.  Also
$\tau=\mathbf1*\mathbf1$.  All convolution rearrangements below are
rearrangements of finite divisor sums.

\begin{lemma}
\label{sb:symmetry}
For $x\ge1$, put
\[
 A(x)=\sum_{n\le x}\bigl(\Lambda(n)\log n+(\Lambda*\Lambda)(n)\bigr).
\]
Then
\begin{equation}
 |A(x)-2x\log x|\le2048x.
 \label{sb:symmetry-bound}
\end{equation}
\end{lemma}

\begin{proof}
Let $L(n)=\log n$ and define $(\mathcal Da)(n)=a(n)\log n$.
Since $\log n=\log d+\log(n/d)$,
\[
 \mathcal D(a*b)=(\mathcal Da)*b+a*(\mathcal Db).
\]
The divisor identity used earlier is $L=\mathbf1*\Lambda$.  Applying
$\mathcal D$ gives the exact pointwise identity
\[
 L^2=L*\Lambda+\mathbf1*(\mathcal D\Lambda)
     =\mathbf1*(\Lambda*\Lambda+\mathcal D\Lambda).
\]
Convolution with $\mu$ therefore gives
\begin{equation}
 \mu*L^2=\Lambda*\Lambda+\mathcal D\Lambda.
 \label{sb:exact-symmetry}
\end{equation}
Similarly $\mathcal D\tau=2L*\mathbf1$.  With the finite parameter
$c$ from \eqref{sb:finite-c} and $f$ from
\eqref{sb:comparison-function},
\begin{equation}
 \mu*f=2L-2c\mathbf1+4c^2\varepsilon.
 \label{sb:comparison-convolution}
\end{equation}

For each $d\le x$, apply \eqref{sb:comparison-error} at $y=x/d$.
Since $|\mu(d)|\le1$, reversing finite sums gives
\begin{align*}
 \left|A(x)-\sum_{n\le x}(\mu*f)(n)\right|
 &\le80\sum_{d\le x}\sqrt{x/d}\bigl(\log(x/d)+1\bigr).
\end{align*}
The preceding integral bound gives
$\log z+1\le4z^{1/4}$ for $z\ge1$, while a finite rectangle
comparison gives
\[
 \sum_{d\le x}d^{-3/4}
       \le1+\int_1^x t^{-3/4}\,dt\le4x^{1/4}.
\]
Consequently the last error is at most $1280x$.

By \eqref{sb:comparison-convolution},
\[
 \sum_{n\le x}(\mu*f)(n)=2P_1(x)-2c\lfloor x\rfloor+4c^2.
\]
Its distance from $2x\log x$ is at most $10x$: indeed
\eqref{sb:log-sum}, $0<c\le1$, and $x\ge1$ bound that distance by
$6x+4\le10x$.  Thus $1290x$ already bounds
$|A(x)-2x\log x|$, and the stated constant $2048$ is valid.
The auxiliary parameter $c$ has disappeared from the conclusion.
\end{proof}

\subsection{The equation for the error term}
\label{sb:error-subsection}

\begin{lemma}
\label{sb:error-equation}
For $x\ge1$, define
\[
 E_R(x)=R(x)\log x+\sum_{n\le x}\Lambda(n)R(x/n).
\]
Then
\begin{equation}
 |E_R(x)|\le4096x.
 \label{sb:error-bound}
\end{equation}
\end{lemma}

\begin{proof}
Integrating the finite step function $\psi$ gives
\[
 \sum_{n\le x}\Lambda(n)\log n
      =\psi(x)\log x-\int_1^x\frac{\psi(t)}t\,dt.
\]
The integral is nonnegative and at most $4x$ by
\eqref{sb:chebyshev-bound}.  Also
\[
 \sum_{n\le x}(\Lambda*\Lambda)(n)
       =\sum_{n\le x}\Lambda(n)\psi(x/n).
\]
The symmetry bound therefore implies
\[
 \left|\psi(x)\log x+\sum_{n\le x}\Lambda(n)\psi(x/n)
                -2x\log x\right|\le2052x.
\]
Substitute $\psi(y)=R(y)+y$.  The difference between the expression
inside this absolute value and $E_R(x)$ is
\[
 x\left(\sum_{n\le x}\frac{\Lambda(n)}n-\log x\right),
\]
whose absolute value is at most $6x$.  Hence
$|E_R(x)|\le2058x\le4096x$.
\end{proof}

\subsection{A second substitution and finite summation by parts}
\label{sb:weighted-sum-subsection}

\begin{lemma}
\label{sb:weighted-sum}
For $x\ge4$,
\begin{equation}
 |R(x)|(\log x)^2
 \le2\sum_{n\le x}\log n\,|R(x/n)|+2^{17}x\log x.
 \label{sb:weighted-sum-bound}
\end{equation}
\end{lemma}

\begin{proof}
Multiply the defining equation for $E_R(x)$ by $\log x$, split
$\log x=\log n+\log(x/n)$ in its sum, and use the same equation at
$x/n$.  This gives the exact identity
\begin{align*}
 R(x)(\log x)^2={}&
   -\sum_{n\le x}\Lambda(n)\log n\,R(x/n)\\
 &+\sum_{ab\le x}\Lambda(a)\Lambda(b)R(x/(ab))\\
 &+E_R(x)\log x-\sum_{n\le x}\Lambda(n)E_R(x/n).
\end{align*}
The last line has absolute value at most
\[
 4096x\left(\log x+\sum_{n\le x}\frac{\Lambda(n)}n\right)
 \le4096x(2\log x+6)
 \le32768x\log x,
\]
because $\log x\ge\log4\ge1$.  The arithmetic function
\[
 a(n)=\Lambda(n)\log n+(\Lambda*\Lambda)(n)
\]
is nonnegative.  Taking absolute values gives
\begin{equation}
 |R(x)|(\log x)^2
    \le\sum_{n\le x}a(n)|R(x/n)|+32768x\log x.
 \label{sb:first-weighted-bound}
\end{equation}

Put $N=\lfloor x\rfloor$, $g_n=|R(x/n)|$ for $1\le n\le N+1$,
and
\[
 \Delta_n=A(n)-2P_1(n).
\]
The symmetry estimate and \eqref{sb:log-sum} give
\[
 |\Delta_n|\le2048n+4n\le4096n.
\]
Since $x/(N+1)<1$, we have $g_{N+1}=x/(N+1)$.  The exact finite
summation-by-parts identity is
\begin{equation}
 \sum_{n=1}^N(a(n)-2\log n)g_n
   =\Delta_Ng_{N+1}
       +\sum_{n=1}^N\Delta_n(g_n-g_{n+1}).
 \label{sb:finite-abel}
\end{equation}
Monotonicity of $\psi$ gives
\[
 |g_n-g_{n+1}|
 \le\psi(x/n)-\psi(x/(n+1))+\frac{x}{n(n+1)}.
\]
In the first term of the resulting weighted sum, each $\Lambda(m)$
is counted with weight $\lfloor x/m\rfloor\le x/m$.  Therefore
\begin{align*}
 \sum_{n=1}^N n|g_n-g_{n+1}|
 &\le x\sum_{m\le x}\frac{\Lambda(m)}m
                   +x\sum_{n=1}^N\frac1{n+1}\\
 &\le x(\log x+6)+x\log(N+1)
 \le x(2\log x+7).
\end{align*}
Here $N+1\le2x$ and $\log2\le1$.  The endpoint term in
\eqref{sb:finite-abel} is at most $4096x$.  Consequently
\[
 \left|\sum_{n\le x}(a(n)-2\log n)|R(x/n)|\right|
 \le4096x(2\log x+8)\le40960x\log x.
\]
Combine this with \eqref{sb:first-weighted-bound}.  The two error
constants add to $73728$, which is less than $2^{17}$.
\end{proof}

\subsection{Replacing the weighted sum by a bounded integral}
\label{sb:sum-integral-subsection}

\begin{lemma}
\label{sb:sum-integral}
For $x\ge1$,
\begin{equation}
 \left|\sum_{n\le x}\log n\,|R(x/n)|
            -\int_1^x\log t\,|R(x/t)|\,dt\right|\le20x.
 \label{sb:sum-integral-bound}
\end{equation}
\end{lemma}

\begin{proof}
Set $N=\lfloor x\rfloor$ and, for $t\ge1$, put
$g(t)=|R(x/t)|$ and $h(t)=\log t\,g(t)$.  The bound
$g(t)\le4x/t$ also holds when $t>x$, because then $R(x/t)=-x/t$.
For $n\le t\le n+1$, monotonicity of $\psi$ gives
\begin{align*}
 |h(t)-h(n)|
 \le{}&\log n\left(\psi(x/n)-\psi(x/t)+\frac xn-\frac xt\right)\\
      &+(\log t-\log n)g(t).
\end{align*}
Integrate over each of the $N$ unit intervals and sum.

For the contribution involving $\psi$, each $\Lambda(m)$ can occur
only in the one unit interval containing $x/m$.  Its contribution is
at most $\Lambda(m)\log(x/m)$, because that interval has length at
most $1$ and $n\le x/m$.  The total is thus at most
\[
 \sum_{m\le x}\Lambda(m)\log(x/m)
      =\int_1^x\frac{\psi(u)}u\,du\le4x.
\]
The contribution from $x/n-x/t$ is at most
\[
 x\sum_{n=1}^N\frac{\log n}{n(n+1)}\le x.
\]
To verify the last inequality using only finite sums, telescope:
\begin{align*}
 \sum_{n=1}^N\log n\left(\frac1n-\frac1{n+1}\right)
 &=\sum_{n=2}^N\frac{\log n-\log(n-1)}n
                         -\frac{\log N}{N+1}\\
 &\le\sum_{n=2}^N\frac1{n(n-1)}\le1.
\end{align*}
Finally, $\log t-\log n\le(t-n)/n$ and $g(t)\le4x/t$ bound the
last contribution by
\[
 4x\sum_{n=1}^N\frac1{n^2}\le8x.
\]
Indeed $1/n^2\le1/(n(n-1))$ for $n\ge2$, so the last sum is at
most $2$.  We have proved
\[
 \left|\sum_{n=1}^Nh(n)-\int_1^{N+1}h(t)\,dt\right|\le13x.
\]
On the interval $[x,N+1]$, $h(t)=x\log t/t\le x$, and the length
of the interval is at most $1$.  Removing that final integral gives
the stronger bound $14x$, hence the displayed bound $20x$.
\end{proof}

\begin{theorem}
\label{sb:contraction-inputs}
Write $r(t)=e^{-t}R(e^t)$ for $t\ge0$.  Then
\begin{align}
 |r(t)|&\le4,
 \label{sb:r-bound}\\
 \left|\int_a^b r(v)\,dv\right|&\le20
          &&(0\le a\le b),
 \label{sb:r-signed}\\
 |r(t)|&\le\frac2{t^2}\int_0^t(t-v)|r(v)|\,dv+\frac{2^{18}}t
          &&(t\ge2).
 \label{sb:r-weighted}
\end{align}
For $0\le u\le v$ one also has
\begin{equation}
 r(v)\ge r(u)-4(v-u).
 \label{sb:r-downward}
\end{equation}
Thus, in the contraction lemma, one may use
$B=4$, $C=20$, $S=2^{18}$, and the harmless larger downward constant
$L=5$.
\end{theorem}

\begin{proof}
The first assertion is \eqref{sb:chebyshev-bound}; the second is
\eqref{sb:signed-integral} after the finite change of variables
$x=e^v$.  Lemmas~\ref{sb:weighted-sum} and \ref{sb:sum-integral} give,
for $x\ge4$,
\[
 |R(x)|(\log x)^2
 \le2\int_1^x\log u\,|R(x/u)|\,du
                    +2^{17}x\log x+40x.
\]
Since $\log x\ge1$, the last two terms are at most
$2^{18}x\log x$.  Put $t=\log x$ and $v=\log(x/u)$.  Direct
substitution yields
\[
 \frac1x\int_1^x\log u\,|R(x/u)|\,du
       =\int_0^t(t-v)|r(v)|\,dv.
\]
Divide by $t^2$ to obtain \eqref{sb:r-weighted}; $t\ge2$ ensures
$e^t\ge4$.

Finally, monotonicity of $\psi$ implies
\[
 r(v)+1=\frac{\psi(e^v)}{e^v}
       \ge e^{-(v-u)}\frac{\psi(e^u)}{e^u}
       =e^{-(v-u)}(r(u)+1).
\]
Because $0\le r(u)+1\le4$ and $1-e^{-h}\le h$ for $h\ge0$,
this gives \eqref{sb:r-downward}.
\end{proof}
